%! AP Statistics — Unit 5, Lesson 5.2 — Homework (Answer Key)
\documentclass[10pt]{article}
\usepackage{apstats-article}
\usepackage{apstats-key}   % pulls in -boxes; do NOT also load -boxes
\usepackage{grading}
\gradingassignment{apstats-u05-l02-hw}
\gradingcourse{advanced_statistics}
\newcommand{\TallMath}[1]{$\displaystyle #1\rule[-1.4em]{0pt}{3.2em}$}

\begin{document}

\pageheader{Unit 5, Lesson 5.2}{Homework --- Answer Key}

\vspace{0.15in}

\begin{enumerate}

  \item Women's heights $N(64.0, 2.8)$.
        \begin{enumerate}[label=\alph*.]
          \item \ans{$z = (60-64)/2.8 \approx -1.43$; $P(X<60) = P(Z<-1.43) \approx 0.0764$.}\gradedkey{\gradedblank[id=hw1a, type=short, points=2, answer=0.0764, check=numeric, tolerance=0.003]{0.9in}{0.28in}}
          \item \ans{$z_1 = (62-64)/2.8 \approx -0.71$, $z_2 = (68-64)/2.8 \approx 1.43$; $P(62<X<68) = 0.9236 - 0.2389 \approx 0.6847$.}\gradedkey{\gradedblank[id=hw1b, type=short, points=2, answer=0.6847, check=numeric, tolerance=0.005]{0.9in}{0.28in}}
          \item \ans{$z = \text{invNorm}(0.90) \approx 1.282$; $x = 64 + 1.282(2.8) \approx 67.6$ in.}\gradedkey{$x = $~\gradedblank[id=hw1c, type=short, points=2, answer=67.6, check=numeric, tolerance=0.1]{0.9in}{0.28in}}
          \boxguard[17]
          \item \ansline{Approximately 90\% of women are shorter than 67.6 inches; the 90th percentile height for women is about 67.6 inches.}\gradedkeyline{\gradedblank[id=hw1d, type=rubric, points=1, display={About 90\% of women are shorter than 67.6 inches.}]{\linewidth+6pt}{\gradedlinesheight{1}}}
        \end{enumerate}

  \item Egg mass $N(62, 5)$ g.
        \begin{enumerate}[label=\alph*.]
          \item \ans{$z = (63-62)/5 = 0.20$; $P(X>63) = 1 - 0.5793 \approx 0.4207$. About 42\% of eggs are large.}\gradedkey{\gradedblank[id=hw2a, type=short, points=2, answer=0.4207, check=numeric, tolerance=0.003]{0.9in}{0.28in}}
          \newpage
          \item \ans{Top 15\%: $z = \text{invNorm}(0.85) \approx 1.036$; $x = 62 + 1.036(5) \approx 67.2$ g. Extra-large eggs weigh at least 67.2 g.}\gradedkey{$x = $~\gradedblank[id=hw2b, type=short, points=2, answer=67.2, check=numeric, tolerance=0.1]{0.9in}{0.28in}}
        \end{enumerate}

  \boxguard[22]
  \item Text messages.
        \begin{enumerate}[label=\alph*.]
          \item \ansline{No. The normal model requires a symmetric, bell-shaped distribution. A strongly right-skewed distribution does not satisfy this condition, so the normal model is not appropriate for the individual observations.}\gradedkeyline{\gradedblank[id=hw3a, type=rubric, points=2, display={No: the distribution is strongly right-skewed, not symmetric and bell-shaped.}]{\linewidth+6pt}{\gradedlinesheight{1}}}
          \boxguard[17]
          \item \ansline{By the Central Limit Theorem (Lesson 5.3), if the student takes many random samples and computes sample means, those sample means will be approximately normally distributed for large enough sample sizes --- even though the original distribution is skewed.}\gradedkeyline{\gradedblank[id=hw3b, type=rubric, points=2, display={Use means of large random samples: by the CLT they are approximately normal.}]{\linewidth+6pt}{\gradedlinesheight{1}}}
        \end{enumerate}

\end{enumerate}

\begin{extensionbox}
\textbf{Extension.} $N(200, 35)$; most extreme 4\% (2\% in each tail).\\
\ans{$z = \pm\text{invNorm}(0.02) \approx \pm 2.054$. Lower: $200 - 2.054(35) \approx 128$ mg/dL. Upper: $200 + 2.054(35) \approx 272$ mg/dL. Adults with cholesterol below 128 or above 272 mg/dL require medical monitoring.}
\end{extensionbox}

\begin{reflectionbox}
\textbf{Preview:} Lesson 5.3 introduces the Central Limit Theorem, which explains why
sampling distributions of means are approximately normal --- even for non-normal
populations --- when samples are large enough.
\end{reflectionbox}

\end{document}
